\begin{abstract}
For every sufficiently large fixed leaf bound \(L\), we prove randomized polynomial-time many-one NP-hardness of realizable Collection Minimum Monotone Satisfying Assignment with weight gap \(L^{1-o(1)}\) and NO satisfaction threshold \(o(1)\). The corresponding bounded-advice learning statement follows through Hirahara--Nanashima's transfer. This establishes the \(\epsilon\)=0 strengthening asked in Section 7 of their June 23, 2026 revision. We combine their star-projection compilation with Minzer--Zheng's Grassmann PCP machinery, an altered repetition parameter order, explicit advice-posterior and zoom-out estimates, a maximal-extension threshold ladder, and finite-list completeness repair. Every leaf/advice bound is fixed independently of input length; the polynomial exponent may depend on it. The contribution is the nearlinear exponent in the realizable regime, with earlier realizable hardness explicitly credited below.

\end{abstract}
\noindent\textbf{Keywords:} hardness of approximation; monotone formulas; realizable learning; PCPs.
\paragraph{Status.} Local submission draft based on the archived version 0.1.0. The full Lean formalization is incomplete; no formal certification is claimed. This paper has not been submitted or announced.
\section{Problems, conventions, and results}\label{sec:1}
A monotone formula is an AND/OR formula with positive variable leaves; repeated occurrences count separately. \(\mathcal F[L]\) denotes formulas with at most \(L\) leaves. Constants arising from empty branches may be simplified away. For \(x\) in \(\{0,1\}^N\) and positive rational weights \(w_i\) with \(\sum_i\) \(w_i\)=1, write \(w(x)=\sum_{i:x_i=1}w_i\). An instance of \(\operatorname{Gap}^{\epsilon,\gamma}_\sigma\mathcal F[L]\text{-CMMSA}\) is an indexed nonempty explicit list \(F_1\),...,\(F_M\), those weights, and a rational budget 0\(<\)\(s\)\(\le\)1. The promises are:

\begin{itemize}
\item YES: some \(x\) with \(w(x)\)\(\le\)\(s\) satisfies at least a 1-\(\epsilon\) fraction of the list.
\item NO: every \(x\) with \(w(x)\)\(\le\)\(\sigma\) \(s\) satisfies strictly less than \(a\) \(\gamma\) fraction.
\end{itemize}

Indices count with multiplicity. Weights and budgets have polynomial bit encodings and inverse-polynomial lower bounds in the number of variables. This is the normalized-weight restriction of HN Definitions 1.7 and 4.1. The distributional version DMMSA specifies the formula distribution by a sampling circuit. Our final output is an explicit power-of-two list, also represented by its exact uniform sampler. Realizable means \(\epsilon\)=0.

A randomized polynomial-time many-one reduction maps each input to one output instance, and for each input satisfies the appropriate YES or NO promise with probability at least 2/3. All randomness and output lengths are polynomially bounded. The theorem asserts one such uniform reduction on input instances for each fixed \(L\), with no common polynomial exponent asserted across \(L\). Limits as \(L\) grows compare this family of separately fixed reductions.

\begin{theorem}[Realizable CMMSA]\label{thm:main}There are functions \(\sigma_L\) in the positive integers and \(\gamma_L\) in (0,1), with

\begin{equation}\label{eq:1}
\frac{\log\sigma_L}{\log L}\longrightarrow1,\qquad\gamma_L\longrightarrow0.
\end{equation}

such that, for every sufficiently large fixed integer \(L\), \(\operatorname{Gap}^{0,\gamma_L}_{\sigma_L}\mathcal F[L]\text{-CMMSA}\) is NP-hard under randomized polynomial-time many-one reductions. The output weights and budget can be required to share a polynomial-magnitude common denominator.
\end{theorem}

For completeness we specify the learning model. Fix the universal machine \(U\) used by HN. In their Definition 3.11, a distribution on \(\{0,1\}^n\) has advice complexity at most \(a\) if a function \(\alpha\) mapping random strings of length \(2^{2n}\) to strings of length at most \(a\) makes \(U(\alpha(r);r,n)\) halt on every \(r\), with failure-output probability at most 1/4 and conditional output statistical distance at most \(2^{-2n}\) from the distribution. The advice can depend on \(r\) and need not be computable. This parameter is not the length of a conventional fixed sampler description. For labeled examples use the total encoded sample length in this definition.

\(\operatorname{Gap}^{\epsilon,\gamma}_\sigma\operatorname{Learn}[a]\) takes unary length and program-budget parameters \(1^n,1^s\) and an example oracle \(E\) on \(\{0,1\}^n\) times \(\{0,1\}\), promised to have sampling advice complexity at most \(a\). Its YES case has a program of description length at most \(s\) running in linear time and agreeing with the labels with probability at least 1-\(\epsilon\). Its NO case bounds the agreement of every program of description length at most \(\sigma\) \(s\) by 1/2+\(\gamma\)/2. The NO assertion does not impose an efficient running-time restriction on those programs. We use HN's machine and encoding convention; the universal overhead below is the associated constant \(c_U\). Our reduction outputs an explicit sampling circuit, so the hardness also holds with that circuit supplied in place of only oracle access.

\begin{corollary}[Realizable learning]\label{cor:learn}For every sufficiently large fixed integer advice cap \(a\), put

\begin{equation}\label{eq:2}
\begin{aligned} L&=a-2\lceil\log_2(a+1)\rceil-c_U,\\ \sigma_a^{\mathrm{learn}}&=\lfloor0.49\sigma_L\rfloor, &\gamma_a^{\mathrm{learn}}&=5\gamma_L.\end{aligned}
\end{equation}

Then \(\operatorname{Gap}^{0,\gamma_a^{\mathrm{learn}}}_{\sigma_a^{\mathrm{learn}}}\operatorname{Learn}[a]\) is NP-hard under randomized polynomial-time many-one reductions. For large \(a\) the parameters are in their required domains, \(\sigma^{\mathrm{\mathrm{learn}}}_a\)=\(a^{1-o(1)}\), and \(\gamma^{\mathrm{\mathrm{learn}}}_a\)=\(o(1)\). In YES instances the linear-time program has agreement exactly one; in NO instances all programs within the expanded description budget have agreement at most 1/2+\(\gamma^{\mathrm{\mathrm{learn}}}_a\)/2.
\end{corollary}

These are approximation-hardness results under randomized reductions. They give no complexity-class separation, one-way-function construction, linear-factor hardness, or growing-parameter reduction. No assertion about general improper PAC learning or circuit size is made.

\section{Relation to previous work}\label{sec:2}
HN revision 1~\cite{HN} proves the nearlinear exponent with vanishing nonzero YES error and asks its realizable strengthening separately from its superconstant-parameter question. Hirahara's FOCS 2022 work~\cite{Hir22} already proves realizable learning hardness; HN explicitly records the corresponding realizable CMMSA factor \(L^\alpha\) for a constant \(\alpha\)\(>\)0, with threshold 1/\(L^\alpha\). The advance here is the nearlinear exponent, not the first realizable hardness theorem. The source contracts are the corrected June 23, 2026 revision, not the withdrawn superconstant regime of an earlier version. \cite{HN,Hir22}.

The MZ~\cite{MZ} higher-query construction gives soundness nearly \(R^{-m}\) with \(m\)+1 queries but its displayed main theorem chooses the alphabet after both positive errors. Our proof changes the construction's parameter order; substituting an inverse-alphabet completeness error into that main theorem alone would be circular. MZ24's~\cite{MZ24} two-query result, even granting its alphabet-before-completeness abstract formulation, does not provide this higher-arity contract: independent two-query repetition spends two queries per repetition. Likewise a perfect-completeness dictatorship test is not by itself this unconditional star-projection PCP. The targeted literature assessment found no exact subsumer, but does not certify novelty or priority. \cite{MZ,MZ24,BKM}.

Slack variables, AND products, concentration, weight rounding, secret sharing, and the star compilation are established techniques. The new argument is their quantitative composition with the modified repetition, actual advice posterior, conditional rank/zoom-out control, and descending threshold ladder. We import the local Grassmann theorems and the outer hardness machinery rather than reprove them. The STOC 2026 full text of MZ was not accessible in the assessment; claims about its numerical details are confined to the accessible arXiv v1, also linked by a coauthor.

\section{Notation and imported contracts}\label{sec:3}
All vector spaces are over \(\mathbb F_2\). \(\operatorname{Grass}(E,d)\) is the set of \(d\)-dimensional linear subspaces of \(E\), with uniform distribution unless stated otherwise. \(\genfrac{[}{]}{0pt}{}{n}{a}_2\) is the Gaussian binomial coefficient counting \(a\)-subspaces of an \(n\)-space. \(\operatorname{SD}\) is total variation distance (one half the L1 distance). \(\operatorname{Zoom}\)[\(Q\),\(W\)] is the uniform \(d\)-subspace \(L\) with \(Q\subseteq L\subseteq W\), where \(d\)=2\(h\). The symbol \(L\) for a subspace is local to the Grassmann argument; the scalar \(L\) in Theorem~\ref{thm:main} always denotes the final leaf bound. Similarly, \(a\) denotes advice dimension in that argument and the final advice cap only in Corollary~\ref{cor:learn}. The final asymptotics use base-two \(\log\), and \(\ln\) is natural \(\log\). The parameters \(m\),\(h\),\(\rho\),\(J\) are leaf-query count, inner dimension parameter, rational decoding slack, and repetition count. Set \(r\)=10\(m\)/\(\rho\), choosing rational \(\rho\) with this value integral. Thus \(r\) is a fixed integer bound on advice dimensions and zoom-out codimensions; the local decoder gives \(a\)+\(c\)\(\le\)\(r\).

The proof uses the following established inputs at their stated parameter ranges. Precise public versions and bibliographic links are in the bibliography. The contracts are explicit mathematical imports; none is a locally formalized theorem.

\paragraph{Imported contract.} \textbf{Outer hardness and game (MZ Theorem 3.1, Section 3.2, Claim 3.2).} There is an absolute NO value \(s_0\)\(<\)1 for 3-Lin, while its fixed positive YES error may be arbitrarily small. Use the bounded-occurrence form with no pair of variables in more than one equation. In the smooth repeated game choose \(J\) equations independently; in each triple keep all coordinates with probability 1-\(\beta\), otherwise a uniform single coordinate. Both provers receive the same \(r\) random vectors on the retained coordinates, extended by zero to the first question. They answer assignments; acceptance requires agreement on retained variables and satisfaction of the \(J\) equations. With fixed NO gap \(\eta\)=1-\(s_0\), its value is at most \(2^{-\Omega(\eta^2 2^{-r} \beta J)}\); its honest failure is at most \(J\) times the outer YES error. In particular the soundness constant does not depend on that later YES error.

\paragraph{Imported contract.} \textbf{Star construction and transport (MZ Section 3.3, Lemmas 3.3--3.4).} Keep tuples \(U\) of \(J\) disjoint equations with no cross-equation pair of variables occurring together in any equation. View \(U\) as a 3\(J\)-space and let \(H_U\) be the span of its equation indicator vectors, of dimension \(J\). Labels at \(L\)+\(H_U\), with \(\dim\) \(L\)=2\(h\) and \(L\) transverse to \(H_U\), are linear functions satisfying the equation right-hand sides on \(H_U\); there are \(R\)=\(2^{2h}\) labels. Center labels are linear functions on a transverse 2(1-\(\rho\))\(h\)-subspace \(K\). Two leaf vertices are equivalent when \(L\)+\(H_U\)+\(H_U'\)=\(L'\)+\(H_U\)+\(H_U'\). The imported lemmas give an equivalence relation and unique side-condition-preserving \(label\) transport. The test chooses \(U\) uniformly, \(K\) uniformly transverse, \(m\) independent transverse \(L_i\) containing \(K\), then a uniform representative from each leaf's equivalence class. It transports each queried \(label\) back and tests restriction to \(K\) against the center \(label\). This is a weighted star-projection (\(m\)+1)-CSP, with alphabets at most \(R\); all choices are finite and enumerable in \(N_{\mathrm{outer}}^{O_m(J)}\). Clique-consistency selection and the collision estimate from Section 5.4.1 preserve the value, up to a vanishing Gaussian-binomial collision loss bounded above by our explicit estimates.

\paragraph{Imported contract.} \textbf{Local decoder (MZ Theorem 4.2).} For fixed small \(\rho\)\(>\)0 and sufficiently large integral \(h\) with integral test dimensions, transverse tables obeying the above side conditions and test density at least \(2^{-2(1-1000\rho)hm}\) have dimensions a,\(c\) with \(a\)+\(c\)\(\le\)10\(m\)/\(\rho\) such that at least \(2^{-6h^2}\) of the uniform \(a\)-subspaces \(Q\) admit \(W\) containing \(Q+H_U\), of codimension \(c\), and a linear \(g\) on \(W\) respecting \(H_U\), with agreement at least \(C\)=\(2^{-2(1-1000\rho^2)h}\)/5 on the conditioned Grassmann space. This is the source construction's scoped contract. For the changed ambient parameter we derive a robust application with input density at least eight times this threshold below, rather than importing an all-\(J\) theorem. Transversality losses are charged explicitly below.

\paragraph{Imported contract.} \textbf{Maximal-pair counting (MZ Definition 5.4; MZ24 revision 1 Theorem 5.26).} For a total table of linear functions on \(\operatorname{Grass}(V,2h)\), a pair (\(W\),\(g\)) is (\(B\),1/5)-maximal relative to \(Q\) if it agrees with the table on at least \(B\) of \(\operatorname{Zoom}\)[\(Q\),\(W\)], and no proper extension \(W'\) with a linear extension of \(g\) agrees on at least \(B\)/5 of \(\operatorname{Zoom}\)[\(Q\),\(W'\)]. Specialize MZ24 with field \(\mathbb F_2\) and its \(\delta\)=\(\rho\)/\(m\), fixing its subsidiary constants as in Lemma 5.24 before taking \(h\) large. For \(\dim\) \(V\)\(\ge\)\(2^h\), \(\dim\) \(Q\)\(\le\)10\(m\)/\(\rho\), \(\operatorname{codim}\) \(W\)\(\le\)10\(m\)/\(\rho\), and \(B\)\(\ge\)\(2^{-2(1-(\rho/m)^3)h}\), the number of these maximal pairs is at most \(B^{-2}\) \(2^{O_{m,\rho}(h)}\). We use the stronger cutoff \(B\)\(\ge\)\(2^{-2(1-\rho^3)h}\), which implies this one. Thus we require no advice range proportional to \(J\) from the broader printed MZ Theorem 5.5. Counting does not assert extension at a single lowered threshold; the threshold ladder below supplies that step.

\paragraph{Imported contract.} \textbf{Covering (KMS~\cite{KMS}, Definition 4.5, Lemmas 4.6--4.7).} For the independent triple-deletion sampler with \(2^d\) \(\beta\)\(\le\)1/8, basic \(d\)-subspace distance is at most \(\beta\) \(\sqrt{J}\) \(2^{d+4}\). For advice dimension \(a\)\(<\)\(d\), except for at most \(\sqrt{\beta}\) \(J^{1/4}\) of uniform \(a\)-subspaces \(Q\), the distance after conditioning on containing \(Q\) is at most \(\sqrt{\beta}\) \(J^{1/4}\) \(2^{d+5}\). The condition is on the joint sampler, changing the distribution of \(V\). These covering lemmas are unconditional results proved in KMS Section 8; they do not require its historical combinatorial hypothesis or UGC.

\paragraph{Imported contract.} \textbf{Compilation (HN Lemma 4.6).} For a star constraint \(e\)=(\(y\);\(x_1\),...,\(x_m\)), let \(\pi_{e,i}\) map its leaf alphabet to the center. Its acceptance is \(\pi_{e,i}(label(x_i))\)=\(label(y)\) for every \(i\). Give a vertex \(u\) occurrence weight \(\lambda(u)=\frac{\mathbb E_e[\mathbf1[y=u]+\sum_i\mathbf1[x_i=u]]}{m+1}\), discarding zero-occurrence vertices. Set \(\Lambda=\sum_u\lambda(u)|\Sigma_u|\), \(w(z_{u,b})\)=\(\lambda(u)\)/\(\Lambda\), and \(s\)=1/\(\Lambda\). If \(X_e\) is the set of distinct leaf vertices, let \(P_{e,x,b}\) contain labels at \(x\) whose projections at every occurrence of \(x\) equal \(b\). Output the monotone formula

\begin{equation}\label{eq:3}
F_e=\bigvee_{b\in\Sigma_y}\left[z_{y,b}\land\bigwedge_{x\in X_e}\left(\bigvee_{a\in P_{e,x,b}}z_{x,a}\right)\right].
\end{equation}

Empty ORs are false. Repeated leaf consistency is imposed in \(P\), not discarded. The formula has at most (\(m\)+1)\(R\) leaves. Value at least 1-\(\tau\) gives a weight-\(s\) assignment satisfying at least 1-\(\tau\) of these formulas. CSP value at most \(\zeta\) implies every assignment of weight at most \(\frac18(5/8)^{1/(m+1)}\) \(\zeta^{-1/(m+1)}\) \(s\) satisfies at most 3/4 of them.

\paragraph{Imported contract.} \textbf{Secret sharing and learning (HN Theorem 3.5 and Lemma 5.1).} Formula leaf bound \(L\) supplies a polynomial-time secret-sharing scheme with total share length at most \(L\). For a circuit-sampled distribution of such formulas, positive inverse-polynomial rational weights and budget, and parameters \(\epsilon\) in [0,1], \(\gamma\) in (0,1], \(\sigma\)\(\ge\)1, the randomized transfer preserves \(\epsilon\), changes \(\gamma\) to 5\(\gamma\) and \(\sigma\) to 0.49\(\sigma\), and gives advice at most \(L\)+2\(\lceil \log_2(L+1)\rceil\)+\(c_U\). Its YES guarantee holds for every transfer random string; its NO guarantee fails with only negligible probability. A sufficiently large polynomial length parameter must make all weighted string lengths integral and satisfy the description/soundness bounds. Our common-denominator step supplies this requirement. Integer gap factors are obtained by rounding down.

We now prove the parameter extension and the remaining composition in full. Lemma~\ref{lem:repair} gives the finite-list repair; Sections~\ref{sec:6}--\ref{sec:10} supply the parameter and reduction arguments for Theorem~\ref{thm:main} and Corollary~\ref{cor:learn}. The local input theorems above retain their usual external-proof status.

\section{Explicit-list exception lemma}\label{sec:4}
\begin{lemma}[Finite-list completeness repair]\label{lem:repair} Input: \(M\)\(\ge\)1 monotone formulas \(F_1\),...,\(F_M\) on \(N\) variables, each with at most \(L\) leaves; positive rational weights \(w_i\) summing to one; positive rational budget \(s\)\(\le\)1; integer \(\sigma\)\(\ge\)8; rational 0\(<\)\(\Gamma\)\(<\)1/2; and rational \(\epsilon\) with \(\epsilon\)\(\ge\)0 and \(\epsilon\) \(\sigma\) \(\le\) \(\Gamma\)/2. The promises are:

\begin{itemize}
\item YES: some assignment of weight at most \(s\) satisfies at least 1-\(\epsilon\) of the list, counting indices with multiplicity.
\item NO: every assignment of weight at most \(\sigma\) \(s\) satisfies strictly less than \(\Gamma\) of the list.
\end{itemize}

Set \(\lambda\)=\(\sigma\) \(s\)/\(\Gamma\). Introduce \(M\) fresh variables \(e_j\) and replace \(F_j\) by \(F_j\) OR \(e_j\). Give original variables raw weights \(w_i\) and every \(e_j\) raw weight \(\lambda\)/\(M\); divide every weight by 1+\(\lambda\). Set

\begin{equation}\label{eq:4}
t=\frac{s+\lambda\epsilon}{1+\lambda},\qquad \sigma_{\mathrm{new}}=\lfloor\sigma/4\rfloor,\qquad \Gamma_{\mathrm{new}}=2\Gamma.
\end{equation}

The resulting list has at most \(L+1\) leaves per formula. In the YES case an assignment of normalized weight at most \(t\) satisfies every formula. In the NO case every assignment of normalized weight at most \(\sigma_{\mathrm{new}}t\) satisfies strictly less than \(\Gamma_{\mathrm{new}}\) of the list.\end{lemma}\begin{proof} In the YES case, take an original witness and enable exactly the exception bits for its failed list indices. Its raw weight is at most \(s\)+\(\lambda\) \(\epsilon\), and it satisfies every augmented formula. No witness or failed-index set is needed to compute the reduction: these choices establish existence.

In the NO case, any augmented assignment of normalized weight at most \(\sigma_{\mathrm{new}}\) \(t\) has raw weight at most

\begin{equation}\label{eq:5}
\lfloor\sigma/4\rfloor(s+\lambda\epsilon)\le\frac{3\sigma s}{8}.
\end{equation}

Its original-variable part therefore satisfies less than \(\Gamma\) of the original list. Its enabled exception bits occupy at most

\begin{equation}\label{eq:6}
\frac{3\sigma s/8}{\lambda}=\frac{3\Gamma}{8}
\end{equation}

of list indices. The augmented acceptance fraction is strictly less than 11 \(\Gamma\)/8, hence strictly less than 2 \(\Gamma\). This proves the claimed perfect-completeness gap reduction. Normalization cancels in both budget comparisons, and \(t\)\(\le\)1 because \(s\)\(\le\)1 and \(\epsilon\)\(\le\)1.

The output has \(N\)+\(M\) variables, \(M\) formulas and at most \(L\)+1 leaves per formula. Construction copies the list, adds \(M\) OR gates, and performs rational arithmetic on input rationals. It is polynomial in the actual input and output bit lengths. For fixed \(\sigma\),\(\Gamma\), inverse-polynomial lower bounds on \(w_i\) and \(s\) imply the same for the normalized weights, \(t\), and \(\lambda\)/\(M\). No SAT, optimum, or counting oracle is used. Applying this directly to an exponential implicit distribution would be invalid; sampling precedes it.\end{proof}

\section{Why the repetition parameters are changed}\label{sec:5}
The alphabet of the MZ star construction is \(R\)=\(2^{2h}\). Its displayed choice \(J\)=\(2^{100h^2}\), \(\beta\)=\(\log_2\log_2 J\)/\(J\) has \(\beta\) \(J\)=\(O(\log h)\), so its outer-game estimate alone does not give an exponent proportional to \(h^2\). We use \(J\)=\(2^{2^{A h^2}}\), \(\beta\)=\(A\) \(h^2\)/\(J\). This observation concerns only arXiv 2510.23991v1; it is not a disproof of its main theorem or a claim about the uninspected STOC 2026 text. The argument below proves the needed parameter order using the local imported contracts.

\section{Parameter extension: posterior and zoom-out argument}\label{sec:6}
This section supplies the parameter modification rather than treating an inverse-alphabet completeness error as part of a black-box theorem. All logarithms here are base two (\(\ln\) denotes the natural logarithm). Fix an integer \(m\)\(\ge\)2 and a positive rational \(\rho\)\(\le\)1/4000. Use the local decoding and maximal-pair counting contracts stated above (MZ Theorem 4.2 and the narrower MZ24 Theorem 5.26 specialization). Their parameters \(r\) and the constants in their bounds depend on \(m\),\(\rho\). The following lemmas justify the changed ambient application with an explicit constant input margin. Choose \(h\) sufficiently large relative to these fixed parameters and a multiple of \(b_m\), the denominator of the selected rational \(\rho\). Then 2(1-\(\rho\))\(h\) and all test dimensions are integers. All subspaces below are over \(\mathbb F_2\), \(U\) consists of \(J\) disjoint triples, and \(\dim(U)\)=3\(J\). The side-condition space \(H_U\) has dimension \(J\).

\begin{lemma}[Exact complement identity for star queries]\label{lem:complement}
Let $U=H\oplus Z$, with $\dim H=J$, $\dim Z=2J$, and $t\le d\le2J$.
Draw a uniform $t$-space $K$ transverse to $H$, then independent uniform
transverse $d$-spaces $L_i\supseteq K$, for $1\le i\le m$.
The law of $(K,H+L_1,\ldots,H+L_m)$ equals the law obtained by drawing
a uniform complement $A$ of $H$, a uniform $t$-space $K\subseteq A$,
and independent uniform $d$-spaces $L_i'\subseteq A$ containing $K$,
and outputting $(K,H+L_1',\ldots,H+L_m')$.
\end{lemma}
\begin{proof}
Fix $K$. Projection to $Z$ expresses it as the graph of a map on a
$t$-space $\overline K$. Each possible $D=H+L$ corresponds to a
$d$-space of $Z$ containing $\overline K$. Its transverse lifts containing
$K$ are exactly the $2^{J(d-t)}$ extensions of that graph map.
Consequently the $D_i$ are independent and uniform among these possible
spaces. Independently choose a uniform complement $A$ containing $K$.
The map $D_i\mapsto A\cap D_i$ is a bijection to the $d$-spaces of $A$
containing $K$, and $H+(A\cap D_i)=D_i$.
There are $2^{J(2J-t)}$ complements containing each transverse $K$,
while every complement contains the same Gaussian-binomial number of
$t$-spaces. Thus the uniform incidence law of $(K,A)$ can be sampled
in either order, proving the assertion.

The actual leaf table is indexed by $D_i=H+L_i$ and its transported
label is restricted to $K$. The acceptance event is therefore preserved
exactly. Random clique representatives can be coupled at identical
$D_i$, with identical transport; alternatively use the already selected
clique-consistent table. No condition on intersections among projected
leaf increments is needed. This is not an identity of the unobserved
raw tuples $(K,L_1,\ldots,L_m)$.
\end{proof}

\begin{lemma}[Inverse agreement with explicit ambient bounds]\label{lem:inverse-explicit}
Fix $m\ge2$, $0<\rho\le1/4000$, and integral $r=10m/\rho$.
Let $h$ be sufficiently large with integral dimensions
$t_0=2(1-\rho)h$, $s_0=2\rho h$, $d=2h$; write
\[
 S=2^{-2(1-1000\rho)hm},\qquad
 e=2^{-2(1-1000\rho^2)h},\qquad b=s_0-1.
\]
For tables of linear functions on the $d$-spaces and $t_0$-spaces
of an $n$-space, suppose the actual test (uniform center and independent
uniform containing leaves) has density $\epsilon\ge S$.
Under the ambient lower bounds below, some nonempty
$\operatorname{Zoom}[Q,W]$ with
$\dim Q+\operatorname{codim}W=r$ admits a linear function with
agreement greater than $e$ with the leaf table.
The constants and the lower cutoffs on $h$ depend only on $m,\rho$.
\end{lemma}
\begin{proof}
We use precisely the following analytic inputs from MZ~\cite{MZ}:
Lemma~4.5 lifts an $(r,e)$-pseudorandom Grassmann indicator to an
$(r,2e)$-pseudorandom basis-invariant matrix indicator;
Theorem~4.6 bounds its level-$i$ norm at every dyadic $p\ge4$ by
$2^{500i^2p}(2e)^{(p-2)/p}$ for $i\le r$; and Lemma~4.7 gives
\[
 \|\mathcal T F^{=i}\|_2^2
 \le \bigl(2^{-i(2\rho h-1)}+3\cdot2^{i-n}\bigr)
       \|F^{=i}\|_2^2.
\]
Their bilinear and global-hypercontractive ancestry remains imported.
We reconstruct the intervening moment, selection and exponent steps
rather than invoke the printed inverse proof without these bounds.

Require $h\ge r$, $e\le1/2$, and
\[
 (2h+2\rho mh)2^{2h-n}\le\tfrac12,\qquad
 b\ge1,\qquad n\ge2h+rb+\log_2 6.
\]
The first condition bounds the probability of any column-rank failure
in the common-center matrix sampler. Conditioned on full rank, the
subspace law is the actual uniform star law; the matrix indicators
vanish otherwise. Thus the star density for sets is at most twice
the corresponding matrix expectation.
For an $(r,e)$-pseudorandom leaf set with lift $F$, decompose
$\mathcal T F=L+H_{\mathrm{high}}$ into levels at most $r$ and above $r$.
Orthogonality and the spectral input give
\[
 \|H_{\mathrm{high}}\|_2^2
 \le2^{-(r+1)b}+3\cdot2^{2h-n}\le2^{-rb}.
\]
Take $\eta=2^{-(2/3)r\rho h}$. Markov's inequality gives
$\Pr[|H_{\mathrm{high}}|>\eta]\le2^r2^{-(2/3)r\rho h}$.
On its complement use
$|L+H_{\mathrm{high}}|^m\le2^{m-1}(|L|^m+\eta^m)$;
on the exceptional event use $0\le\mathcal TF\le1$.

Choose a power of two
$T\ge\max(4,4(m+3)/(m\rho))$, and let $P$ be the least power of two
at least $mT$. These depend only on $m,\rho$.
For the lifted center indicator $G$, whose expectation is at most its
Grassmann density $\beta$, H\"older contributes $\beta^{1-1/T}$.
Probability-space norm monotonicity, contraction of $\mathcal T$,
and the dyadic norm input at $P$ imply
\[
 \|L\|_{mT}^m
 \le(r+1)^m2^{500mr^2P}(2e)^{m-2m/P}
 \le(r+1)^m2^{500mr^2P+m}e^{m-2/T}.
\]
The last inequality uses $P\ge mT$ and $2e\le1$.
Combining the splitting and rank factors, the actual matching-star
density $X$ for these two sets satisfies
\[
 X\le K_{\mathrm{inv}}\beta^{1-1/T}e^{m-2/T}
       +D_{\mathrm{inv}}2^{-(20/3)mh},
\]
where the deliberately conservative constants are
\[
 K_{\mathrm{inv}}=2^{2m}(r+1)^m2^{500mr^2P},\qquad
 D_{\mathrm{inv}}=2^m+2^{r+1}.
\]
In particular we do not use the stronger printed high-degree error
or a norm theorem at a possibly nondyadic exponent $mT$.

For the selection step, a good tuple has its $m$ quotient increments
$L_i/R$ in joint direct sum, not merely pairwise trivial intersection.
Its total span has dimension $t_0+ms_0$, so labels agreeing on the
center glue to a linear function on that span. Expose ordered leaf
bases modulo the center. At increment index $j$, the previously
exposed span has at most $2^{t_0+j}$ vectors and the allowed sampling
denominator is at least $2^{n-1}$ when $n\ge d+1$. A union bound gives
\[
 q=\Pr[\text{not jointly direct}]<2^{t_0+ms_0+1-n}\le S/2
 \quad\text{if } n\ge t_0+ms_0+2+\log_2(1/S).
\]
Thus good accepting tuples have mass at least $\epsilon/2$.
This joint-span condition is needed for gluing arbitrary labels in the
inverse argument; it is not a loss in Lemma~\ref{lem:complement}'s
exact quotient-query identity.

Choose a uniform global linear $f$. Let $X_f$ be the unrestricted
star density for which all leaf and center labels match $f$, and let
$\beta_f$ be its matching-center density. Then
\[
 \mathbb E X_f\ge(\epsilon/2)2^{-(t_0+ms_0)},\qquad
 \mathbb E\beta_f=2^{-t_0}.
\]
Average the affine expression
$X_f-(\epsilon/4)2^{-ms_0}\beta_f$ to choose $f$ such that, writing
$X=X_f$ and $\beta=\beta_f$,
\[
 X\ge(\epsilon/4)2^{-ms_0}\beta
          +(\epsilon/4)2^{-(t_0+ms_0)},\qquad
 \beta\ge X\ge(\epsilon/4)2^{-(t_0+ms_0)}
                \ge\tfrac14 2^{-2(m+1)h}.
\]
These are densities over all actual stars, not densities renormalized
by the generic tuples.

Suppose this $f$'s matching leaf set were $(r,e)$-pseudorandom.
To absorb the analytic error against half of the first selected
summand, it suffices to require
\[
 D_{\mathrm{inv}}2^{-(20/3)mh}
 \le (S^2/32)2^{-(t_0+2ms_0)}.
\]
Equivalently, the following fixed-parameter lower cutoff suffices:
\[
 \bigl((8/3)m-2+(3996m+2)\rho\bigr)h
       \ge\log_2(32D_{\mathrm{inv}}).
\]
Its coefficient is positive. This compares with the full selected
signal, including $\beta$, rather than with $S$ alone.
The selected lower bound and analytic upper bound would now imply
\[
 1\le8K_{\mathrm{inv}}2^{ms_0}\beta^{-1/T}
                  e^{m-2/T}/\epsilon
 \le8K_{\mathrm{inv}}2^{2/T}2^{-\lambda h},
\]
where
\[
 \lambda=2000m(\rho-\rho^2)-2m\rho-2(m+3)/T
 \ge m\rho\bigl(2000(1-\rho)-2.5\bigr)>0.
\]
Taking $h>(\log_2(8K_{\mathrm{inv}})+2/T)/\lambda$ is a contradiction.
By MZ Definition~2.1, failure of $(r,e)$-pseudorandomness provides
$\dim Q+\operatorname{codim}W=r$ and matching density greater than
$e$ on a nonempty zoom. The restriction $f|_W$ supplies the asserted
agreement. This uses the definition's actual dimensions and gives
agreement $>e$, not an unspecified constant times $e$.
All constants above are independent of $n$. Every ambient requirement
is an explicit lower bound, satisfied by $n=2J$ for the prescribed
$J$ and sufficiently large admissible $h$; there is no ambient upper
bound or required identity involving $\log\log n$.
\end{proof}

\begin{lemma}[Robust enlarged-ambient local application]\label{lem:robust-local}
Fix $m\ge2$, $0<\rho\le1/4000$ and $r=10m/\rho$, with integral test dimensions
$t=2(1-\rho)h$, $d=2h$. Put
\[
 S=2^{-2(1-1000\rho)hm},\qquad
 \epsilon'=2^{-2(1-1000\rho^2)h},\qquad C=\epsilon'/5.
\]
For sufficiently large admissible $h$ and $J=2^{2^{Ah^2}}$, a transverse
star test of density at least $8S$ in $U$, with $\dim U=3J$ and
$\dim H=J$, has fixed $a,c$ with $a+c\le r$ such that at least
$2^{-6h^2}$ of all uniform $a$-spaces $Q\subseteq U$ admit a
codimension-$c$ space $W\supseteq Q+H$ and a linear function respecting
$H$ with agreement at least $C$ on the transverse conditioned
Grassmann space. Here $A>0$ is fixed before $h$.
\end{lemma}
\begin{proof}
Apply Lemma~\ref{lem:inverse-explicit} at $n=2J$, retaining all its
fixed-parameter cutoffs on $h$ and explicit lower bounds on $n$.
It supplies agreement greater than $\epsilon'=e$ at input density
at least $S$. We now give the amplification and side-condition steps.

By Lemma~\ref{lem:complement}, the mean density over complements is
at least $8S$. Complements whose density $\epsilon_A$ is at least
$4S$ have mass at least $4S$, since the mean is at most $4S$ plus
that mass. Fix such a complement. Starting from its table, repeatedly
apply the inverse argument while the current density is at least
$\epsilon_A/2$. Add the resulting zoom to a union $X$ of leaf entries,
and independently replace every entry in $X$ by a uniform linear
function. The inverse threshold remains valid since
$\epsilon_A/2\ge2S$.

Here is the preservation event, including its ambient bound.
For each candidate $(Q,W',g)$ with dimensions bounded by $r$, either
$X$ has relative measure less than $2^{-2h}$ in the zoom, or the
fraction of the entries in $X$ agreeing with $g$ is at most
$2^{1-2h}$. Each refreshed entry agrees with a fixed $g$ with probability
$2^{-2h}$. The Chernoff bound and the lower bound
$|X\cap\operatorname{Zoom}[Q,W']|
 \ge2^{-2h}2^{(2h-r)(n-r-2h)}$
in the second case give a per-step failure probability at most
\[
 F_n=16(r+1)^2 2^{n(r+1)}
       \exp\!\left(-2^E/12\right),\qquad
 E=-4h+2+(2h-r)(n-r-2h).
\]
Indeed at most $16(r+1)^2 2^{n(r+1)}$ triples are needed: the Gaussian
counts for $Q,W'$ contribute at most $16\cdot2^{nr}$ and the number
of linear $g$ is at most $2^n$. The same bound holds conditionally on
any preceding history, since each refresh is independent.

Require $2^{1-2h}\le\epsilon'/2$. Before a new zoom is added, at most
this much of its agreement comes from previously refreshed entries;
this includes the alternative with relative measure below $2^{-2h}$.
Its unrefreshed fraction is therefore at least $\epsilon'/2$.
The zoom has ambient measure at least $2^{-rn}/4$ by the elementary
Gaussian bounds, so $X$ increases by at least
$\epsilon'2^{-rn}/8$. Thus no more than
$B_n=\lceil8\cdot2^{rn}/\epsilon'\rceil+1$ such steps are possible.
Choose the parameters so $B_nF_n<1/2$; a history with all preservation
events then exists. This requirement is satisfied for the prescribed
$J$: when $h\ge\max(r+4,16)$ and $n\ge8h+4r+16$, we have
$E\ge hn/2$, whereas $\log B_n+\log(16(r+1)^2 2^{n(r+1)})$
is $O_{m,\rho}(rn+h)$. Every selected function agrees with the original
table on at least $\epsilon'/2$ of its zoom.

At termination the acceptance has fallen by at least $\epsilon_A/2$.
It can change only if at least one of the $m$ leaf occurrences is in
$X$. Each leaf marginal is uniform, including when leaves repeat,
so $\mu(X)\ge\epsilon_A/(2m)$. For a fixed $a$-space $Q$, the
family of all $d$-spaces containing $Q$ has measure at most
$2^{4h^2-an}$. Count distinct $Q$ rather than algorithmic visits:
a repeated $Q$ with another zoom still covers only this family.
For each distinct useful $Q$, choose one witnessing $W'$ and assign
its codimension $c$. Let $N_{a,c}$ count the assigned pairs with
$\dim Q=a$, let $T_a$ be the number of all $a$-spaces in the complement,
and put $f_{a,c}=N_{a,c}/T_a$. Since $T_a\le4\cdot2^{an}$,
\[
 \frac{\epsilon_A}{2m}\le\mu(X)
 \le\sum_{a,c}N_{a,c}2^{4h^2-an}
 \le4\cdot2^{4h^2}\sum_{a,c}f_{a,c}.
\]
Thus $\sum_{a,c}f_{a,c}\ge\epsilon_A2^{-4h^2}/(8m)$.
Set these fractions to zero on bad complements and average over all
uniform complements. Good complements have mass at least $4S$ and
$\epsilon_A\ge4S$, so the sum of the averaged fractions is at least
$2S^2 2^{-4h^2}/m$. A single pigeonhole over at most $(r+1)^2$
dimension pairs gives one fixed $(a,c)$ whose success probability is
at least $2S^2 2^{-4h^2}/(m(r+1)^2)$, and hence at least the
conservative bound retained below:
\[
 \frac{2S^2}{m(r+1)^3}2^{-4h^2}.
\]
The pair is fixed after pigeonholing. We do not condition the marginal
law on a successful complement: successes form a subset of the full
uniform experiment. Its $Q$ marginal is uniform among $H$-disjoint
$a$-spaces. Extending to uniform $Q\subseteq U$ loses at most
$2^{a+1-2J}$. The sufficient bounds
\[
 h^2\ge4mh+\log_2(m(r+1)^3)+1,\qquad 2J\ge5h^2+r+2
\]
leave at least $2^{-5h^2-1}\ge2^{-6h^2}$ of all $Q$.
Finally set $W=W'\oplus H$ and extend $g$ by the prescribed side
conditions. For fixed $A,Q,W'$, projection to $W/H$ gives the same
uniform quotient law for the containing leaves in $W$ and $W'$;
the values on $H$ already agree. Hence their agreement events agree
and the retained $\epsilon'/2$ is stronger than $C=\epsilon'/5$.
This accounts for the uniform-$Q$ correction inside this lemma;
subsequent exclusions are applied to its already obtained lucky set.
\end{proof}

Write \(a\) for an advice dimension, 0\(\le\)\(a\)\(\le\)\(r\), and \(c\) for a zoom-out codimension, 0\(\le\)\(c\)\(\le\)\(r\). Write \(d\)=2\(h\). In the outer sampler each triple retains all three coordinates with probability 1-\(\beta\), and otherwise retains a uniformly chosen single coordinate. Thus the number of dropped blocks \(D\) is \(Bin(J,\beta)\), and \(\dim(V)\)=3\(J\)-2\(D\). Let \(P\) be the uniform distribution on \(a\)-subspaces \(Q\) of \(U\). Let \(P'\) be the marginal obtained by first sampling \(V\), then sampling a uniform \(a\)-subspace of \(V\).

For \(J\)=\(2^{2^{A h^2}}\) and \(\beta\)=\(A\) \(h^2\)/\(J\), the advice statistical distance obeys

\begin{equation}\label{eq:7}
\operatorname{SD}(P,P')\le\delta_{\mathrm{adv}}:=\beta\sqrt J\,2^{a+4}.
\end{equation}

This is smaller than \(2^{-100h^2}\) for sufficiently large \(h\). Likewise the KMS zoom-in exceptional fraction

\begin{equation}\label{eq:8}
\delta_{\mathrm{zoom}}:=\sqrt\beta\,J^{1/4}
\end{equation}

and the conditional distance \(\delta_{\mathrm{zoom}}\) \(2^{d+5}\) are smaller than \(2^{-100h^2}\), after increasing the fixed lower bound on \(h\). The condition \(2^d\) \(\beta\)\(\le\)1/8 is satisfied. Increasing \(A\) does not cause a circular choice: \(A\) depends only on fixed decoding/outer constants, and \(h\) is chosen afterward.

Set \(T\)=\(h^4\) and \(\zeta\)=\(2^{-30h^2}\). The binomial tail has the explicit bound

\begin{equation}\label{eq:9}
\Pr[D>T]\le\left(\frac{eAh^2}{T}\right)^T\le2^{-100h^2}
\end{equation}

for sufficiently large \(h\) (round \(T\) to an integer if needed). Except on a \(P'\)-mass at most \(2^{-70h^2}\) of \(Q\), Markov's inequality gives \(\Pr[D>T\mid Q]\)\(\le\)\(\zeta\). Passing from \(P'\) to \(P\) adds at most \(\delta_{\mathrm{adv}}\) to this exceptional mass. Further, the set where \(P'(Q)\)\(<\)\(P(Q)\)/2 has \(P\)-mass at most 2 \(\delta_{\mathrm{adv}}\), by the definition of statistical distance. Intersect these good sets with the KMS good zoom-in set. The excluded \(P\)-mass is at most

\begin{equation}\label{eq:10}
\delta_{\mathrm{zoom}}+3\delta_{\mathrm{adv}}+2^{-70h^2}<2^{-20h^2}.
\end{equation}

Fix a \(Q\) outside this exceptional set. For any \(V\) with \(D\)\(\le\)\(T\), Bayes' rule and the Gaussian binomial formula give

\begin{equation}\label{eq:11}
\begin{aligned}\frac{\Pr[V\mid Q]}{\Pr[V]}&=\frac{\mathbf1[Q\subseteq V]}{\genfrac{[}{]}{0pt}{}{\dim V}{a}_2P'(Q)}\\&\le2\frac{\genfrac{[}{]}{0pt}{}{3J}{a}_2}{\genfrac{[}{]}{0pt}{}{3J-2D}{a}_2}\le8\cdot2^{2aT}.\end{aligned}
\end{equation}

The last inequality follows by factoring out \(2^{2aD}\) in the product formula; the remaining finite products are bounded by an absolute constant because 3\(J\)-2\(D\) is much larger than \(a\). This is a density bound for the actual posterior, not an assumption that \(V\) remains unconditioned.

Now let \(W\)=\(W(Q)\) be any subspace containing \(Q\), of codimension \(c\)\(\le\)\(r\). Once \(Q\) is fixed, \(W\) is fixed before \(V\) is sampled. Express \(W\) as the kernel of a rank-\(c\) matrix. Failure of \(\operatorname{codim}_V(W \cap  V)\)=\(c\) means that some nonzero linear combination of its rows has support entirely outside \(V\). For each of the at most \(2^c\)-1 combinations, select one nonzero coordinate. That coordinate is removed with probability at most \(\beta\). A union bound therefore gives the unconditional estimate (\(2^c\)-1)\(\beta\), uniformly in \(W\). Applying the posterior density bound on \(D\)\(\le\)\(T\), and then the tail bound, gives, for this same \(Q\),

\begin{equation}\label{eq:12}
\begin{aligned}\Pr[\operatorname{codim}_V(W\cap V)\ne c\mid Q]&\le8\cdot2^{2aT}(2^c-1)\beta+\zeta\\&\le2\zeta.\end{aligned}
\end{equation}

The last inequality holds for large \(h\) because \(J\) is double exponential in \(h^2\). The estimate is valid for every \(W(Q)\) of the permitted codimension; there is no union over all \(W\), no assumption that \(W\) is chosen independently of \(Q\), and no use of an unconditional rank estimate after conditioning.

The KMS distribution conditioned on \(Q\) is exactly the posterior needed here. Indeed, the probability \(a\) uniform \(d\)-subspace \(L\) of \(V\) contains \(Q\) is

\begin{equation}\label{eq:13}
\mathbf1[Q\subseteq V]\frac{\genfrac{[}{]}{0pt}{}{d}{a}_2}{\genfrac{[}{]}{0pt}{}{\dim V}{a}_2}.
\end{equation}

The numerator is independent of \(V\), so conditioning the joint (\(V\),\(L\)) sampler on \(Q\) subset \(L\) induces the same \(V\) posterior as sampling \(Q\) uniformly inside \(V\). This also explains why an unconditioned outer draw is invalid.

For the zoom-out, put \(b\)=\(d\)-\(a\) and \(p_0\)=\(2^{-cb}\). On the rank-stable event, the additional probability of \(L\) subset \(W\), given \(V\) and \(Q\), is

\begin{equation}\label{eq:14}
p_V=\frac{\genfrac{[}{]}{0pt}{}{\dim V-a-c}{b}_2}{\genfrac{[}{]}{0pt}{}{\dim V-a}{b}_2}=p_0\bigl(1+O(2^{-J/2})\bigr).
\end{equation}

This product estimate is uniform: \(\dim(V)\)\(\ge\)\(J\) always, while \(a\),\(c\),\(d\) are \(O_{m,\rho}(h)\). Under the uniform \(U\) sampler the corresponding probability is at least \(p_0\)/2 for large \(h\). Conditioning two distributions at distance \(\Delta\) on an event of probability at least \(p_0\)/2 changes their distance by at most \(O\)(\(\Delta\)/\(p_0\)). Apply this with the KMS conditional distance \(\Delta\)=\(\delta_{\mathrm{zoom}}\) \(2^{d+5}\).

Finally compare the posterior mixture with the mixture additionally weighted by \(p_V\). Rank-unstable \(V\) have posterior mass at most 2 \(\zeta\); their contribution to the normalized weighted mixture is \(O\)(\(\zeta\)/\(p_0\)). On the remaining \(V\), \(p_V\)/\(p_0\) differs from one by \(O(2^{-J/2})\). Hence for any event concerning \(L\), the difference between these two mixtures is

\begin{equation}\label{eq:15}
O(2^{-J/2}+\zeta/p_0).
\end{equation}

Both this expression and \(\Delta\)/\(p_0\) are much smaller than \(2^{-10h^2}\), since \(p_0\)\(\ge\)\(2^{-2rh}\). Therefore agreement with a decoded function on a uniform \(L\) satisfying \(Q\subseteq L\subseteq W\) transfers to the actual posterior experiment with additive error below \(2^{-10h^2}\). This supplies the conditional-on-\(W\) step with the correct \(V\) distribution. It also supplies the rank stability required when the two provers extend their functions.

\subsection{Decoder probability is measured in the same joint law}
The outer advice consists of shared random vectors in \(V\) (MZ Section 3.2), and both provers take the span of the same first \(a\) vectors after independently guessing \(a\). Given linear independence, the span is uniform among \(a\)-subspaces of \(V\). The probability of dependence is at most \(\sum_{i=0}^{a-1}2^{i-\dim V}\) \(\le\) \(2^{r-J}\). Thus the actual vector-advice law and the ideal experiment \(V\) then uniform \(Q\) subset \(V\) differ in statistical distance by at most \(2^{r-J}\). This error is negligible compared with every success term below; no alteration of the outer game is needed.

Fix the leaf-\(label\) assignment \(T_1\) after the imported clique-consistency selection, with \(T_2\) the center-\(label\) assignment. A good first-prover question \(U\) is one whose conditional inner-test density reaches the robust threshold 8\(S\) in the preceding lemma, where \(S\)=\(2^{-2(1-1000\rho)hm}\). Write \(\Delta\)=\(2^{-2(1-\xi)hm}\). After the already charged clique loss, the surviving density is at least \(\Delta\)/2. If \(\Delta\)/4\(\ge\)8\(S\), the mass of good \(U\) is at least \(\Delta\)/4, since the mean is at most 8\(S\) plus that mass. The condition \(\Delta\)/\(S\)\(\ge\)32 is exactly 2hm(\(\xi\)-1000\(\rho\))\(\ge\)5. It holds for sufficiently large \(h\) with 1000\(\rho\)\(\le\)\(\xi\)/4. Thus the good-question mass remains \(2^{-O_m(h)}\), with the amplification margin paid explicitly and no additional complement-coupling error.

Fix a good first-prover question \(U\). In the ideal joint experiment the marginal law of \(Q\) is \(P'\), and the conditional law of \(V\) is exactly the posterior analyzed above. The favorable advice set has \(P\)-mass at least \(2^{-6h^2}\)-\(2^{-20h^2}\)-\(2^{r+1-2J}\). Passing to \(P'\) subtracts \(\delta_{\mathrm{adv}}\), so its \(P'\)-mass is at least \(2^{-8h^2}\). For each such \(Q\) the expected agreement under the actual posterior \(V\) is at least \(C\)-\(2^{-10h^2}\), hence at least \(C\)/2. Since agreement is in [0,1], the posterior probability that it is at least \(C\)/4 is at least \(C\)/4. Subtracting the rank-failure probability 2 \(\zeta\) leaves at least \(C\)/8 for large \(h\).

These two probabilities multiply within the actual joint experiment: \(P'(good Q)\) times \(\Pr\)[good \(V\) \(\mid\) \(Q\)]. They are never converted to an unconditional \(V\) success probability for a fixed \(Q\). In particular, the factor 8\(\cdot\)\(2^{2a h^4}\) was used solely to upper-bound a rank-failure event; its reciprocal is not a loss in the decoded success bound. The independent-vector correction subtracts at most \(2^{r-J}\) from the final joint success, again negligible. The threshold-ladder strategy below supplies the maximal-extension step.

\subsection{Total tables and transverse subspaces}
The source vertices \(L\) plus \(H_U\) require \(L\cap H_U=\{0\}\). For each \(U\), extend the induced table to all \(d\)-subspaces by assigning the zero linear function on otherwise undefined entries. The robust enlarged-ambient local lemma above is applied with its proved side-condition/transversality hypothesis. Its resulting agreement on the full conditioned Grassmann space, or equivalently the agreement after charging the exceptional entries, loses a negligible amount quantified here.

If \(Q\) intersect \(H_U\)=\{0\} and \(W\) contains \(Q+H_U\) with codimension at most \(r\), a uniform \(d\)-subspace \(L\) between \(Q\) and \(W\) fails transversality only if \(L/Q\) meets (\(H_U\)+\(Q\))/\(Q\) nontrivially. A union bound over nonzero vectors of this \(J\)-dimensional image bounds this probability by

\begin{equation}\label{eq:16}
2^{J+d-\dim W+1}\le2^{-2J+d+r+1}.
\end{equation}

This follows from \(\Pr\)[\(z\) in a uniform \(b\)-subspace of an \(n\)-space] =(\(2^b\)-1)/(\(2^n\)-1), and applies with an extra factor at most two if a preceding distribution was already conditioned on transversality. It is negligible compared with \(C\).

For each second-prover question \(V\) fix an arbitrary legitimate completion \(U_0(V)\) containing \(V\). Define its total table on every \(d\)-subspace \(L\) of \(V\) by \(T_1\)[\(L\)+\(H_{U_0}\)]\(\mid\)\(L\) when transverse, and by zero otherwise. This definition depends only on \(V\) and the fixed global assignment table. Whenever \(L\) is transverse to both \(H_U\) and \(H_{U_0}\), the source clique-consistency relation identifies the two values. No equality is asserted on other entries.

Here are explicit charges for those other entries in the actual joint law for a fixed original \(U\). On \(D\)\(\le\)\(T\), \(\dim(V)\)\(\ge\)3\(J\)-2\(T\). For uniform \(Q\) subset \(V\), the probability \(Q\) meets \(H_{U_0(V)}\) nontrivially is at most \(2^{a+J-\dim(V)+1}\) \(\le\) \(2^{-2J+2T+r+1}\). This bound is valid although \(U_0\) is selected after \(V\). Including the tail gives joint probability at most \(2^{-100h^2}\)+\(2^{-2J+2T+r+1}\). By Markov, except on \(P'\)-mass at most \(2^{-69h^2}\), its posterior probability for fixed \(Q\) is at most \(\zeta\). Passing to \(P\) adds \(\delta_{\mathrm{adv}}\). Remove these \(Q\) from the good advice set; the total removed advice mass still fits the bound \(2^{-20h^2}\).

For remaining \(Q\),\(V\) with \(D\)\(\le\)\(T\), rank stability, and \(Q\) disjoint from both side-condition spaces, take \(L\) uniform between \(Q\) and \(W\) intersect \(V\). For either side-condition space the same quotient union bound gives nontransversality probability at most

\begin{equation}\label{eq:17}
2^{J+d-(\dim V-r)+1}\le2^{-2J+2T+d+r+1}.
\end{equation}

For \(H_{U_0}\) not contained in \(W\) intersect \(V\) use its intersection with that space; its dimension is at most \(J\), which only improves this bound. Thus replacing either total table by the actually defined transverse values costs at most twice this bound on good \(V\), plus the posterior tail and exceptional probabilities \(O(\zeta)\). These are below \(2^{-10h^2}\). Enlarge the finite lower bound on \(h\) to absorb all such additive errors simultaneously. In the subsequent agreement ledger one may use expected agreement \(C\)/2, good-\(V\) agreement \(C\)/4, and good-\(V\) probability \(C\)/8 as already budgeted. No undefined table entry or unproved identity on nontransverse \(L\) is used.

\subsection{Explicit threshold ladder for maximal extensions}
A single arbitrarily lowered threshold does not guarantee maximality. Instead the second prover guesses \(j\) uniformly from \(\{0,\ldots,r\}\) and uses \(B_j\)=\(C\)/(4\(\cdot\)\(5^j\)). For a favorable \(V\), begin with (\(W(Q)\cap V\),\(g\) restricted to \(W(Q)\cap V\)), which has codimension at most \(r\) and agreement at least \(B_0\)=\(C\)/4. At stage \(j\), if the pair is (\(B_j\),1/5)-maximal, stop. Otherwise Definition 5.4 supplies a proper extension of the pair with agreement at least \(B_j\)/5=\(B_{j+1}\). Replace the pair by that extension. Each replacement strictly decreases codimension, so this process stops by stage \(r\), at latest at the whole space, where no proper extension exists.

Thus at one of these \(r\)+1 thresholds a maximal pair exists whose function extends the first prover's restriction. The second prover enumerates all maximal pairs of codimension at most \(r\) at the guessed threshold and chooses uniformly; this is a strategy for a value bound, not a polynomial algorithm used by the reduction. The maximal-pair counting contract bounds each such list by \(B_j^{-2}\)\(\cdot\)\(2^{O_{m,\rho}(h)}\)=\(2^{O_{m,\rho}(h)}\). Its threshold hypothesis holds: the local agreement is \(C\)=\(2^{-2(1-1000\rho^2)h}\)/5, so

\begin{equation}\label{eq:18}
\frac{B_r}{2^{-2(1-\rho^3)h}}=\frac{2^{2(1000\rho^2-\rho^3)h}}{20\cdot5^r}\ge1
\end{equation}

for sufficiently large \(h\) (take \(\rho\)\(<\)1). Codimension is at most the source local bound \(r\)\(\le\)10\(m\)/\(\rho\); the advice dimension is also within the narrower MZ24 counting contract. This costs the additional constant factor 1/(\(r\)+1), with no single-threshold existence assumption. The first prover takes its decoded linear function \(g\) on \(W(Q)\), which respects the equation side conditions, and extends it uniformly to \(U\). The second prover extends its selected maximal-pair function uniformly to \(V\). For the successful choice, that function extends \(g\) restricted to \(W(Q)\cap V\). Rank stability gives \(\operatorname{codim}_V\)(\(W(Q)\cap V\))=\(c\), where \(c\)=\(\operatorname{codim}_U\) \(W(Q)\)\(\le\)\(r\). The restriction to \(V\) of a uniform extension of \(g\) to \(U\) is uniform among the \(2^c\) extensions of this common restriction: the quotient map \(V\)/(\(W(Q)\cap V\)) \(\to\) \(U\)/\(W(Q)\) is an isomorphism. Thus it agrees with the second prover's chosen extension on \(V\) with probability \(2^{-c}\)\(\ge\)\(2^{-r}\). The first prover's assignment satisfies all outer equations because its extension preserves \(g\) on \(H_U\). This proves the claimed random-extension success cost.

\section{Completing the modified PCP parameter order}\label{sec:7}
Fix \(m\)\(\ge\)2 and a desired exponent slack \(\xi\)\(>\)0; set \(\rho\)\(\le\)\(\min(\xi/4000,1/4000)\), choosing it positive and rational and shrinking it further if necessary to meet the imported local theorems. Use the same star-query construction as MZ, with the new \(J\),\(\beta\) above. The following table checks every ambient/repetition-dependent loss used in its decoding argument. Local Grassmann decoding and the maximal-zoom-out counting theorem are imported; the numerical parameter extension is argued here.

\begin{longtable}{@{}p{0.29\textwidth}p{0.67\textwidth}@{}}\toprule
\textbf{Use} & \textbf{Control under the replacement}\\\midrule\endhead
Ambient dimensions in local decoding & Counting uses \(\dim(V)\)\(\ge\)\(J\)\(\ge\)\(2^h\). Complement decoding has ambient \(n\)=2\(J\) and its query law is exact. The robust local lemma retains the matrix-rank, spectral, randomization-history and uniform-\(Q\) bounds; each ambient error improves with \(J\).\\[5pt]
Duplicate clique choices & The source Gaussian-binomial comparison decreases exponentially in \(J\); for large \(h\) the union over \(m^2\) pairs is at most \(2^{-J}\).\\[5pt]
Lucky advice meeting \(H_U\) & Its uniform probability is at most \(2^{r+1-2J}\).\\[5pt]
Advice/zoom-in covering & The bounds above and their exceptional \(Q\) sets are less than \(2^{-20h^2}\), below the local decoding lucky mass \(2^{-6h^2}\).\\[5pt]
\(\operatorname{Zoom}\)-out and rank steps & The posterior calculation above gives error below \(2^{-10h^2}\), below the decoded agreement \(C\)\(\ge\)\(2^{-O_{m,\rho}(h)}\).\\[5pt]
Number of maximal zoom-outs & The imported bound is \(2^{O_{m,\rho}(h)}\) for the relevant agreement and codimension; its constants do not depend on \(J\).\\[5pt]
Random extensions & Their agreement cost is at most \(a\) factor \(2^r\), since codimensions are bounded by \(r\).\\[5pt]
\bottomrule\end{longtable}

Write the good-question mass as at least $K_U^{-1}2^{-B_Uh}$ and the
maximal-list bound as at most $K_M2^{B_Mh}$, with $K_U,K_M\ge1$ and
$B_U,B_M\ge0$. These name the bounds just established and imported;
their constants depend on $m,\rho$, not $A,J$ or the later YES error.
The displayed decoding strategy, with its previous losses already
charged, succeeds before the vector-law correction with probability
at least $K^{-1}2^{-8h^2-Bh}$, where
\[
 K=40K_UK_M(r+1)^3 2^r,\qquad B=B_U+B_M+2.
\]
Here $C/8\ge2^{-2h}/40$, the three dimension/threshold guesses cost
at most $(r+1)^3$, and the extension cost is at most $2^r$.
For $h\ge\max(1,B+\log_2K)$ this is at least $2^{-9h^2}$.
Taking $J\ge9h^2+r+1$, the vector-law error $2^{r-J}$ leaves at
least $2^{-9h^2-1}\ge2^{-10h^2}$. No exponential-in-$J$ success
factor is paid.

The fixed outer NO gap gives success at most \(2^{-\kappa \beta J}\), with \(\kappa\)\(>\)0 depending on the fixed advice size and outer NO gap, not on the later outer YES error. Choose an integer A before \(h\):

\begin{equation}\label{eq:19}
A>\frac{20}{\kappa}.
\end{equation}

Then the outer upper bound, even including the factor two for legitimate
conditioning ensured by the padding below, is at most
$2\cdot2^{-\kappa Ah^2}<2^{-10h^2}$, contradicting decoded success.
All lower bounds on $h$ may depend on this fixed $A$. Thus the modified
PCP has soundness at most $R^{-(1-\xi)m}$, where $R=2^{2h}$.

Only after fixing \(m\),\(\xi\),\(h\),\(J\),\(\beta\) do we fix any desired positive constant PCP completeness error \(\tau\). Select the outer 3-Lin YES error

\begin{equation}\label{eq:20}
\epsilon_1\le\frac{\tau}{100(m+1)J}.
\end{equation}

The starting hardness theorem allows every fixed positive \(\epsilon_1\) with an absolute NO gap. In particular, this does not require exactly satisfiable linear equations to be NP-hard. Its reduction exponent may depend on \(\epsilon_1\), which is allowed because every parameter here is a constant independent of the input length.

To account for discarded illegitimate tuples, take sufficiently many disjoint copies of the outer instance so their probability \(a\) is at most \(\min(\tau/100,1/4)\). The bound is \(O(J^2/N_{\mathrm{outer}})\) with a constant from the bounded occurrence condition. Copies preserve the value and degree bound. The number required depends only on the already fixed parameters; for an input of length \(n\) the padded size is \(O(n+N_{\min})\), still polynomial. Conditioning on legitimate \(U\) multiplies a failure probability by at most 1/(1-\(a\)). The sampler's marginal for each clique-resampled \(U'_i\) is uniform over legitimate \(U\): the initial vertex is uniform, and uniform resampling within its equivalence class preserves that measure. All relevant spaces have the same dimensions and number of transverse extensions.

An honest global labeling satisfies the composed constraint whenever the equations in the original \(U\) and all \(m\) resampled \(U'_i\) are satisfied. Union bounding gives failure at most

\begin{equation}\label{eq:21}
\frac{(m+1)J\epsilon_1}{1-a}\le\frac{\tau}{75}<\tau.
\end{equation}

Using all \(m\)+1 blocks is deliberately conservative; no equality between a single-block bound and the multi-query acceptance event is needed. This establishes the required parameter order: for each fixed \(m\),\(\xi\) and each sufficiently large \(h\), every fixed positive \(\tau\) is available without increasing \(R\). \(J\) and the source input-size threshold can be very large, and all are charged as fixed constants.

The construction and enumeration cost is \(N_{\mathrm{outer}}^{O_{m}(J)}\) with constants depending on these choices, plus the starting outer reduction. It is not a polynomial with a uniform exponent when \(L\), \(h\) or \(\tau\) vary with input length. No superconstant-parameter hardness claim is made.

\section{From the modified PCP to an explicit realizable list}\label{sec:8}
Use HN's star-projection compilation and its occurrence weights, retaining repeated leaf-variable consistency as in its Lemma 4.6. It gives formulas with at most (\(m\)+1)\(R\) leaves, a normalized honest budget \(s\), and a constant NO satisfaction bound 3/4 up to weight \(\sigma\) \(s\), where we may safely take the explicit integer

\begin{equation}\label{eq:22}
\sigma=\left\lfloor\frac{R^{1-1/m}}{16}\right\rfloor,\qquad\xi=\frac1{m^2}.
\end{equation}

For implementation this floor is computed by integer root comparison; no irrational number is an encoded weight. The constant 1/16 is smaller than the coefficient in the source's list-decoding bound. For sufficiently large \(R\), \(\sigma\)\(\ge\)8. Set \(q\)=\(\lfloor \sqrt{m}\rfloor\) and take an AND of \(q\) independently sampled base formulas. Put

\begin{equation}\label{eq:23}
\Gamma=2(3/4)^q,\qquad\epsilon=\frac{\Gamma}{16\sigma},\qquad\tau=\frac{\epsilon}{4q}.
\end{equation}

The modified PCP above permits this \(\tau\) after \(R\) has been fixed. The AND-product YES failure is at most \(q\) \(\tau\)=\(\epsilon\)/4; the NO satisfaction fraction is at most \((3/4)^q\)=\(\Gamma\)/2. The local leaf bound is \(q(m+1)\)\(R\).

Let \(N_0\) be the number of monotone variables. Materialize the finite sampling support if necessary; for fixed parameters it has polynomial size in the outer input. Its rational probabilities and occurrence weights have polynomial bit descriptions. A positive atom is a product of a fixed number of inverse-polynomial choices, so positive variable weights and \(s\) are inverse-polynomially bounded. Adding rational numbers may enlarge \(a\) common denominator, which is handled separately below.

Construct a bounded-random-bit sampler whose distribution has statistical distance at most \(\epsilon\)/8 from the exact product distribution. One implementation enumerates its \(S\) atoms and rounds cumulative rational probabilities to a dyadic grid with at least \(\log_2(8S/\epsilon)\) bits. Charge the support enumeration and rational arithmetic. This is polynomial for fixed \(L\); it is not an oracle for an exponentially large support.

Choose a natural \(P\) with \(1/\epsilon\le P\); \(\epsilon\) is positive. Since \(\epsilon\) is rational, \(P=\lceil 1/\epsilon\rceil\) can be computed by integer division and remainder on its numerator and denominator. Use the conservative count

\begin{equation}\label{eq:24}
\begin{aligned}T&=32(N_0+11)P^2, &M&=2^{\operatorname{clog}_2 T},\\\frac{N_0\ln2+\ln12}{2(\epsilon/8)^2}&\le T\le M<2T.\end{aligned}
\end{equation}

Here \(\operatorname{clog}_2 T\) is the least natural exponent whose power of two is at least \(T\). The constructor uses only natural arithmetic. The displayed analytic bound follows from \(\ln 2\le1\), \(\ln 12\le11\) and \(1/\epsilon\le P\); it is a proved comparison, not a real-number computation required by the sampler. This one count already gives sampling success at least 5/6, and therefore also the required 2/3. It is the least qualifying power of two for \(T\), not necessarily for the exact logarithmic threshold. For fixed \(L\) the previously chosen positive rational \(\epsilon\) is fixed, so this choice of \(P\) is constant with respect to the outer input, and \(M<64(N_0+11)P^2\) bounds the list length. Support enumeration and rational arithmetic costs still have to be charged as above. The formal count and probability bounds do not certify an encoded machine runtime.

Draw \(M\) independent formulas from this sampler. Hoeffding's inequality and a union bound over \(2^{N_0}\) assignments imply, with probability at least 2/3, that all empirical satisfaction fractions differ from their exact product probabilities by at most \(\epsilon\)/4, including the sampling-distribution error. Hence the list's YES failure is at most \(\epsilon\)/2, and its NO acceptance is at most \(\Gamma\)/2+\(\epsilon\)/4 \(<\) \(\Gamma\). The exception lemma may therefore use the conservative error budget \(\epsilon\). Its condition \(\epsilon\) \(\sigma\)\(\le\)\(\Gamma\)/2 holds with a factor-eight margin.

Apply that lemma to this actual list. The output satisfies all constraints in its YES case, with ratio \(\lfloor \sigma/4\rfloor\), threshold 2 \(\Gamma\), and at most \(q(m+1)\)\(R\)+1 leaves per formula. Failure of the randomized reduction occurs only in the sampling event already bounded by 1/3. The completed list has an exact uniform sampler using \(\log_2(M)\) random bits. In particular, the later learning reduction sees \(\epsilon\)=0 exactly; it does not inherit the approximate sampler used to construct the list.

\section{Rational-weight control and learning transfer}\label{sec:9}
Polynomial bit lengths alone do not make all weights integral multiples of 1/poly(\(n\)). We supply a rounding step before any use of fixed-length strings \(w_i\) \(\lambda\) in the learning reduction. Let \(v_i\) be the final normalized weights on \(N'\) variables and \(t\) its YES budget. Choose the least power of two \(D\)\(\ge\)8(\(N'\)+1)/\(t\), set \(a_i\)=\(\lceil D v_i\rceil\), \(A_0\)=\(\sum_i\) \(a_i\), and

\begin{equation}\label{eq:25}
B_0=\min\{A_0,\lceil Dt\rceil+N'\},\qquad v_i'=a_i/A_0,\qquad t'=B_0/A_0.
\end{equation}

A YES assignment remains feasible, since its integer weight is at most \(D\) \(t\)+\(N'\). Conversely, any assignment with \(v'\)-weight at most (\(\sigma_{\mathrm{new}}\)/2)\(t'\) has original \(v\)-weight at most

\begin{equation}\label{eq:26}
\frac{\sigma_{\mathrm{new}}}{2}\frac{B_0}{D}\le\frac{\sigma_{\mathrm{new}}}{2}\left(t+\frac{N'+1}{D}\right)\le\frac9{16}\sigma_{\mathrm{new}}t<\sigma_{\mathrm{new}}t.
\end{equation}

Thus the NO acceptance threshold is unchanged if the gap factor is rounded down to \(\lfloor \sigma_{\mathrm{new}}/2\rfloor\). Formula leaves and satisfaction are unchanged. \(D\) and \(A_0\) are polynomial numeric quantities because \(t\) has an inverse-polynomial lower bound; all new weights and the budget share the polynomial common denominator \(A_0\). The clipping in \(B_0\) guarantees \(t'\)\(\le\)1 and does not weaken either displayed inequality.

Represent the power-of-two list by its exact sampler. Every monotone formula with \(L\) leaves admits the source's \(L\)-bit-total secret-sharing construction. HN Lemma 5.1 explicitly allows \(\epsilon\)=0: it keeps that completeness error, multiplies the NO threshold by five, loses a factor 0.49 in the size gap, and uses advice bound

\begin{equation}\label{eq:27}
L+2\lceil\log_2(L+1)\rceil+c_U.
\end{equation}

Its length parameter can be chosen as a polynomial multiple of \(A_0\), so all strings of length \(v'_i\) times that parameter have integral lengths. The parameter is also increased by a polynomial factor to meet its stated description and soundness bounds. These adjustments are polynomial in the actual instance size. Perfect completeness holds for every randomness string of this transfer; its additional randomized failure probability is on the NO side. The conservative list count above already gives success at least 5/6; combine this with the transfer's sufficiently small failure probability to obtain a randomized many-one reduction with success at least 2/3. Finitely many small inputs can be handled by a fixed truth table, a constant cost depending on the fixed parameters.

\section{Final asymptotic parameter choice}\label{sec:10}
For each fixed sufficiently large target leaf bound \(L\), consider integers \(m\)\(\ge\)256 with \(m\)\(\le\)\(\sqrt{\log_2 L}\). Let \(q\)=\(\lfloor \sqrt{m}\rfloor\) and

\begin{equation}\label{eq:28}
h(L,m)=b_m\left\lfloor\frac{\log_2((L-1)/(q(m+1)))}{2b_m}\right\rfloor,\qquad R=2^{2h}.
\end{equation}

Additionally require \(b_m\)\(\le\)\(\sqrt{\log_2 L}\). Require \(h(L,m)\) to exceed every fixed lower bound used above for that \(m\) and \(\xi\)=1/\(m^2\), including \(\sigma\)\(\ge\)8 and all numerical error inequalities. Choose the largest admissible \(m\). Every fixed \(m\) is admissible for all sufficiently large \(L\), so \(m(L)\) tends to infinity. The finite selection does not assert a polynomial-time algorithm uniform in \(L\); for each fixed \(L\) it fixes the constants of a uniform polynomial-time reduction on inputs. The same fixed-parameter convention is used in the revised source.

The final leaf bound is at most \(L\) by construction. Moreover,

\begin{equation}\label{eq:29}
\begin{aligned}\log R&=\log L-O(\log(m+1)+b_m)=\log L-o(\log L),\\\frac{\log\sigma}{\log L}&\longrightarrow1,\\\Gamma&=2(3/4)^{\lfloor\sqrt m\rfloor}\longrightarrow0.\end{aligned}
\end{equation}

All floor operations, the exception lemma, weight rounding and learning transfer lose only constant factors in \(\sigma\). They therefore preserve \(\sigma\)=\(L^{1-o(1)}\). For the learning target with advice cap \(a\), choose \(L=a-2\lceil\log_2(a+1)\rceil-c_U\); for large \(a\) this is positive and the advice overhead fits within \(a\). Also \(L\)/\(a\) tends to one, preserving the exponent.

Theorem~\ref{thm:main} follows with the explicit parameters

\begin{equation}\label{eq:30}
\sigma_L=\left\lfloor\frac{\lfloor\sigma/4\rfloor}{2}\right\rfloor,\qquad\gamma_L=2\Gamma.
\end{equation}

where \(\sigma\), \(\Gamma\), \(m\), \(h\) and \(R\) are selected above. For sufficiently large \(L\) these are in the promised ranges. The conservative natural list count above gives sampling success at least 5/6. Corollary~\ref{cor:learn} follows from the stated HN transfer with its residual error made at most 1/6; the combined success is at least 2/3. The exact sampler and complete repaired YES witness ensure that the learning error parameter is exactly zero. This completes the proofs of Theorem~\ref{thm:main} and Corollary~\ref{cor:learn}.\hfill$\square$

\section{Research and review disclosure}\label{sec:11}
This manuscript was developed through AI-assisted mathematical derivation, source challenge, independent agent mathematical review, dependency audit, and source/significance assessment under Quantyra Research. These are informal AI reviews, not independent human peer review or Lean verification. The source challenger contributed to the argument and is not counted as an independent mathematical verifier. Review roles, findings, and limits are disclosed in the companion review disclosures. Bibliographic versions and imported inputs are recorded in the bibliography.
